% COSC428 Computer Vision - Conference Paper Template
% Project: CNN-Based Sun Elevation Detection for Image Geolocation
% Based on IEEE conference paper format (two-column)
%
% HOW TO COMPILE:
%   pdflatex cosc428_paper_template.tex
%   bibtex   cosc428_paper_template
%   pdflatex cosc428_paper_template.tex
%   pdflatex cosc428_paper_template.tex
%
% INSTRUCTIONS TO AUTHOR:
%   - Replace all [PLACEHOLDER] blocks with your actual content.
%   - Every \todo{} comment is a reminder of what belongs in that location.
%   - Target: 6 full pages. Adjust figure sizes / text density to fill space.
%   - Remove all \todo{} calls and this comment block before submission.

\documentclass[conference]{IEEEtran}

% -----------------------------------------------------------------------
% Packages
% -----------------------------------------------------------------------
\usepackage{cite}
\usepackage{amsmath,amssymb,amsfonts}
\usepackage{graphicx}
\usepackage{textcomp}
\usepackage{xcolor}
\usepackage{hyperref}
\usepackage{booktabs}       % nicer tables
\usepackage{subcaption}     % subfigures
\usepackage{algorithm}
\usepackage{algpseudocode}

% Helper command – highlights things you still need to write.
% Remove or redefine before submission.
\newcommand{\todo}[1]{\textcolor{red}{\textbf{[TODO: #1]}}}

% -----------------------------------------------------------------------
% Title & Authors
% -----------------------------------------------------------------------
\title{Solar Elevation Estimation from Single Images Using Convolutional Neural
Networks for Multi-Image Geolocation}

\author{
    \IEEEauthorblockN{[Alex Nowlan]}
    \IEEEauthorblockA{
        Department of Computer Science and Software Engineering \\
        University of Canterbury, Christchurch, New Zealand \\
        \texttt{[ano58@uclive.ac.nz]}
    }
}

\begin{document}

\maketitle

% -----------------------------------------------------------------------
% ABSTRACT  (~150–200 words)
% -----------------------------------------------------------------------
% PURPOSE: Give readers a self-contained snapshot of the entire paper.
%          Write this LAST, after all other sections are finalised.
%
% MUST COVER:
%   1. The problem being solved (what + why it matters).
%   2. Your approach at a high level (CNN fine-tuned from ImageNet, sun
%      position model, triangulation pipeline).
%   3. Key experimental result(s) with concrete numbers.
%   4. The main conclusion / takeaway.
\begin{abstract}
\todo{Write the abstract last. Cover: (1) problem statement – estimating
geographic location from consumer photographs without GPS metadata;
(2) your approach – fine-tuning an ImageNet-pretrained CNN (ResNet/VGG/etc.)
to regress solar elevation angle from a single image, combined with an
astronomical sun-position model to infer latitude/longitude from a set of
temporally spaced, co-located images; (3) key quantitative result, e.g.\
``median location error of X\,km on our test set''; (4) significance –
passive, metadata-free geolocation applicable to archival or forensic
photography.}
\end{abstract}

\begin{IEEEkeywords}
solar elevation, geolocation, convolutional neural network, transfer learning,
sun position model, ImageNet, image forensics
\end{IEEEkeywords}

% -----------------------------------------------------------------------
% 1. INTRODUCTION  (~0.5–0.75 page)
% -----------------------------------------------------------------------
% PURPOSE: Motivate the problem, set context, and state contributions.
%
% STRUCTURE (one paragraph each):
%   P1 – Hook: why geolocating images matters (journalism, forensics,
%        wildlife monitoring, autonomous systems, etc.).
%   P2 – Existing approaches and their limitations (EXIF GPS, landmark
%        recognition, scene matching – what they can't do).
%   P3 – The solar-elevation insight: the sun's apparent position is a
%        deterministic function of time, date and observer coordinates;
%        a CNN can learn to estimate elevation from visual cues (shadows,
%        sky brightness, haze). Introduce your pipeline at a high level.
%   P4 – Contributions (use a short bulleted list if you prefer):
%         • A fine-tuned CNN for solar elevation regression.
%         • A closed-form geolocation algorithm using ≥2 timestamped images.
%         • Quantitative evaluation on [your dataset name / source].
%   P5 – Paper organisation sentence: "The remainder of this paper is
%        structured as follows: …"
\section{Introduction}

\todo{Write the introduction following the five-paragraph structure above.
Be specific: name the CNN backbone you chose and why (e.g.\ ResNet-50 –
good accuracy/speed trade-off, widely benchmarked). Mention that
timestamps must be available in EXIF or supplied manually.}

The automatic determination of where a photograph was taken — without
relying on embedded GPS metadata — is an open problem in computer vision
with applications in digital forensics~\cite{placeholder_forensics},
investigative journalism, and autonomous navigation.
\todo{Expand this opening paragraph and add real citations.}

% Example of how to cite: \cite{key}
% Add references in the bibliography section at the end.

% -----------------------------------------------------------------------
% 2. RELATED WORK  (~0.5–0.75 page)
% -----------------------------------------------------------------------
% PURPOSE: Situate your work in the literature. Show you know what has
%          been tried before and how your approach differs.
%
% SUBSECTIONS TO CONSIDER:
%   2A – Image-based geolocation (IM2GPS, PlaNet, Translocator, etc.)
%   2B – Solar position estimation from images (shadow geometry,
%        sky appearance models)
%   2C – Transfer learning / fine-tuning CNNs for regression tasks
%   2D – Astronomical sun-position algorithms (e.g. NREL SPA)
%
% TIP: Do NOT just summarise papers. Explain how each relates to / differs
%      from your approach. End the section with a sentence such as
%      "Unlike prior work, our method requires no landmarks, maps, or
%       reference databases – only the image and its timestamp."
\section{Related Work}

\subsection{Image-Based Geolocation}
\todo{Discuss landmark / scene-matching approaches (e.g.\ IM2GPS~\cite{placeholder_im2gps},
PlaNet~\cite{placeholder_planet}). Note they require large retrieval databases
and fail on featureless outdoor scenes.}

\subsection{Solar Cues in Computer Vision}
\todo{Cover work that uses shadows or sky appearance to estimate sun direction
(e.g.\ hold-geoffroy2017 deep sky model, lalonde2012 outdoor illumination).
Contrast with your regression-based, end-to-end CNN approach.}

\subsection{Transfer Learning for Regression}
\todo{Briefly explain why fine-tuning an ImageNet-pretrained network for a
regression target (elevation angle in degrees) is sound practice. Reference
a general transfer learning paper and one domain-specific example.}

\subsection{Astronomical Sun-Position Models}
\todo{Cite the algorithm you use (e.g.\ NREL Solar Position Algorithm,
Grena 2012, or pysolar). Note accuracy (sub-0.01° for SPA).}

% -----------------------------------------------------------------------
% 3. METHODOLOGY  (~1.5–2 pages – the core technical section)
% -----------------------------------------------------------------------
% PURPOSE: Describe your system so clearly that someone else could
%          reproduce it. Use equations and a system diagram.
\section{Methodology}

% ---- 3A: System Overview ----
% Include a pipeline figure here – e.g. a block diagram showing:
%   [Input Image] --> [CNN Backbone] --> [Elevation Angle θ]
%   Multiple (image, timestamp) pairs --> [Sun Model] --> [Lat/Lon estimate]
\subsection{System Overview}

\todo{Insert your pipeline diagram (Fig.~\ref{fig:pipeline}) and write
2–3 sentences describing the end-to-end flow.}

\begin{figure}[t]
    \centering
    % Replace with your actual figure file:
    % \includegraphics[width=\linewidth]{figures/pipeline.pdf}
    \fbox{\rule{0pt}{3cm}\rule{\linewidth}{0pt}} % placeholder box
    \caption{System pipeline. A CNN estimates solar elevation $\hat{\theta}$
    from each input image. Given timestamps, the sun-position model inverts
    the elevation to candidate (latitude, longitude) loci. Multiple images
    constrain the solution to a unique geographic location.}
    \label{fig:pipeline}
\end{figure}

% ---- 3B: CNN Architecture & Transfer Learning ----
\subsection{CNN Architecture and Transfer Learning}

\todo{State which backbone you used (e.g.\ ResNet-50 pretrained on
ImageNet-1K~\cite{placeholder_imagenet}). Describe the modification:
replace the final classification head with a single-neuron regression
head outputting elevation $\theta \in [-90°, 90°]$. Discuss:
\begin{itemize}
  \item Layer freezing strategy (which layers were frozen / fine-tuned).
  \item Loss function (MSE or MAE for regression).
  \item Activation on output neuron (none, or tanh scaled to $[-90,90]$).
\end{itemize}}

The backbone network $f_\phi$ maps an input image $\mathbf{I} \in \mathbb{R}^{H \times W \times 3}$
to a predicted solar elevation angle:
\begin{equation}
    \hat{\theta} = f_\phi(\mathbf{I}), \quad \hat{\theta} \in [-90°,\, 90°].
    \label{eq:cnn}
\end{equation}

\todo{Explain the training objective, e.g.\
$\mathcal{L} = \frac{1}{N}\sum_{i=1}^{N}(\hat{\theta}_i - \theta_i)^2$.}

% ---- 3C: Dataset and Ground-Truth Labels ----
\subsection{Dataset and Ground-Truth Labels}

\todo{Describe:
\begin{itemize}
  \item Where your training images came from (Flickr, personal collection,
        public dataset). State the total number and train/val/test split.
  \item How ground-truth elevation angles were computed from EXIF timestamps
        and known GPS coordinates using your sun-position library.
  \item Any data augmentation applied (random crop, colour jitter, etc.).
  \item Class / angle distribution (histogram would be great as a figure).
\end{itemize}}

% ---- 3D: Sun-Earth Geometric Model ----
\subsection{Sun--Earth Geometric Model}

\todo{Explain the astronomical model you use to compute solar elevation $\theta$
given observer latitude $\phi$, longitude $\lambda$, date, and UTC time $t$:
\begin{equation}
    \theta = g(\phi, \lambda, t),
    \label{eq:sunmodel}
\end{equation}
where $g(\cdot)$ is the Solar Position Algorithm (or equivalent). Cite the
algorithm. Note the precision achievable ($<0.01°$ for SPA).}

% ---- 3E: Geolocation from Multiple Images ----
\subsection{Geolocation via Multi-Image Constraint Solving}

\todo{This is the novel contribution of your pipeline. Explain:
\begin{itemize}
  \item For a single image with timestamp $t_k$ and predicted elevation
        $\hat{\theta}_k$, equation~\eqref{eq:sunmodel} yields a locus of
        $(\phi, \lambda)$ pairs (a curve on the Earth's surface).
  \item Two or more such curves (from different times) intersect at the
        true location.
  \item How you solve for the intersection (e.g.\ least-squares
        optimisation, grid search, or Levenberg--Marquardt on the residual
        $\sum_k (g(\phi,\lambda,t_k)-\hat{\theta}_k)^2$).
  \item Any ambiguities (e.g.\ hemisphere ambiguity) and how you handle them.
\end{itemize}
A small diagram showing the intersecting loci on a map would strengthen this
sub-section.}

The geolocation estimate $(\hat{\phi}, \hat{\lambda})$ minimises:
\begin{equation}
    (\hat{\phi}, \hat{\lambda}) = \arg\min_{\phi,\lambda}
    \sum_{k=1}^{K} \bigl(g(\phi, \lambda, t_k) - \hat{\theta}_k\bigr)^2.
    \label{eq:optim}
\end{equation}

% ---- 3F: Implementation Details ----
\subsection{Implementation Details}

\todo{List hyperparameters and software stack concisely (a table is clean):
backbone, input resolution, batch size, optimiser (Adam, SGD), learning rate
schedule, number of epochs, hardware (GPU model), framework (PyTorch /
TensorFlow), sun-position library (pysolar, pvlib, or custom). This lets
readers reproduce your work.}

\begin{table}[h]
\centering
\caption{Implementation hyperparameters}
\label{tab:hyperparams}
\begin{tabular}{ll}
\toprule
\textbf{Parameter} & \textbf{Value} \\
\midrule
Backbone            & ResNet-50 (ImageNet-1K pretrained) \\
Input resolution    & $224 \times 224$ \\
Batch size          & \todo{X} \\
Optimiser           & Adam ($\beta_1{=}0.9$, $\beta_2{=}0.999$) \\
Initial LR          & \todo{X} \\
LR schedule         & \todo{e.g.\ cosine decay} \\
Epochs              & \todo{X} \\
Sun-position lib    & \todo{e.g.\ pysolar 0.9} \\
Framework           & \todo{PyTorch X.X} \\
Hardware            & \todo{e.g.\ NVIDIA GTX 1080 Ti} \\
\bottomrule
\end{tabular}
\end{table}

% -----------------------------------------------------------------------
% 4. EXPERIMENTS AND RESULTS  (~1.5–2 pages)
% -----------------------------------------------------------------------
% PURPOSE: Demonstrate that your system works and quantify how well.
\section{Experiments and Results}

% ---- 4A: Evaluation Metrics ----
\subsection{Evaluation Metrics}

\todo{Define all metrics you report:
\begin{itemize}
  \item \textbf{Elevation MAE / RMSE} (degrees): primary CNN accuracy metric.
  \item \textbf{Geolocation error} (km): great-circle distance between
        predicted and ground-truth coordinates.  State whether you report
        mean, median, or top-X\% accuracy.
\end{itemize}}

% ---- 4B: CNN Elevation Estimation Results ----
\subsection{CNN Solar Elevation Estimation}

\todo{Report test-set MAE and RMSE. Compare with:
\begin{itemize}
  \item A naive baseline (e.g.\ predict mean training elevation for all images).
  \item Other backbone choices if you ablated them.
\end{itemize}
Include a scatter plot of predicted vs.\ ground-truth elevation angles
(Fig.~\ref{fig:scatter}) and a residual distribution histogram.}

\begin{figure}[t]
    \centering
    \begin{subfigure}[b]{0.48\linewidth}
        \centering
        \fbox{\rule{0pt}{2.5cm}\rule{\linewidth}{0pt}}
        \caption{Predicted vs.\ ground-truth elevation.}
        \label{fig:scatter}
    \end{subfigure}
    \hfill
    \begin{subfigure}[b]{0.48\linewidth}
        \centering
        \fbox{\rule{0pt}{2.5cm}\rule{\linewidth}{0pt}}
        \caption{Residual histogram.}
        \label{fig:residual}
    \end{subfigure}
    \caption{CNN elevation estimation performance on the test set.
    \todo{Fill in actual MAE/RMSE numbers in the caption or text.}}
    \label{fig:elevation_results}
\end{figure}

\begin{table}[h]
\centering
\caption{Elevation estimation accuracy (test set)}
\label{tab:elevation_accuracy}
\begin{tabular}{lcc}
\toprule
\textbf{Model} & \textbf{MAE (°)} & \textbf{RMSE (°)} \\
\midrule
Mean baseline           & \todo{} & \todo{} \\
\todo{Backbone A}       & \todo{} & \todo{} \\
\todo{Backbone B}       & \todo{} & \todo{} \\
\textbf{Ours (final)}   & \todo{} & \todo{} \\
\bottomrule
\end{tabular}
\end{table}

% ---- 4C: Geolocation Results ----
\subsection{Geolocation Accuracy}

\todo{Report geolocation results. Suggested structure:
\begin{itemize}
  \item Table or figure showing median / mean error (km) as a function of
        the number of input images $K$ (e.g.\ $K = 2, 3, 4, 6$).
  \item A world map or local map showing predicted vs.\ true locations for
        representative test sequences (Fig.~\ref{fig:map}).
  \item Discuss upper-bound performance using ground-truth elevations instead
        of CNN estimates, to separate CNN error from model error.
\end{itemize}}

\begin{figure}[t]
    \centering
    \fbox{\rule{0pt}{4cm}\rule{\linewidth}{0pt}} % replace with map figure
    \caption{Example geolocation results.  \textcolor{green}{$\star$} = ground
    truth; \textcolor{red}{$\bullet$} = predicted location.
    \todo{Replace placeholder with an actual map figure.}}
    \label{fig:map}
\end{figure}

% ---- 4D: Ablation Study ----
\subsection{Ablation Study}

\todo{Optional but strongly recommended. Test the effect of:
\begin{itemize}
  \item Number of images $K$ on geolocation error.
  \item Degree of fine-tuning (head only vs.\ last block vs.\ full network).
  \item Input resolution.
  \item Time-of-day distribution in training data.
\end{itemize}}

% ---- 4E: Qualitative Analysis ----
\subsection{Qualitative Analysis}

\todo{Show 2–3 interesting examples (success cases and failure cases).
For each, display the image, the predicted elevation, and the resulting
location estimate. Discuss what visual cues the model appears to be using
(shadow length/direction, sky colour gradient, haze) using a Grad-CAM or
similar visualisation if possible.}

% -----------------------------------------------------------------------
% 5. DISCUSSION  (~0.5 page)
% -----------------------------------------------------------------------
% PURPOSE: Interpret results, explain failures, and connect back to the
%          research questions posed in the introduction.
\section{Discussion}

\todo{Address the following points:
\begin{itemize}
  \item \textbf{What worked well?} Which conditions (clear sky, strong
        shadows, golden hour) produce the most accurate elevation estimates?
  \item \textbf{What failed and why?} Overcast skies, indoor/artificial
        lighting, extreme latitudes, near-solstice dates (low sun-movement
        rate), timestamp errors.
  \item \textbf{Error budget:} How much of the final geolocation error is
        attributable to CNN elevation error vs.\ inherent geometric ambiguity?
  \item \textbf{Practical considerations:} Minimum number of images needed
        for reliable geolocation; effect of image spacing in time; hemisphere
        ambiguity.
  \item \textbf{Comparison to GPS accuracy:} Put your error (km) in context.
\end{itemize}}

% -----------------------------------------------------------------------
% 6. CONCLUSION AND FUTURE WORK  (~0.3–0.4 page)
% -----------------------------------------------------------------------
% PURPOSE: Summarise contributions and point to what comes next.
\section{Conclusion}

\todo{Write 2–3 paragraphs:
\begin{enumerate}
  \item Restate the problem and your key contributions (1–2 sentences each,
        do not copy the abstract verbatim).
  \item Your best result in one sentence (e.g.\ ``With four co-located images
        our system achieves a median geolocation error of X km without any
        landmark database or prior map knowledge.'').
  \item Future work: fine-tuning on more diverse data, adding moon/star cues
        for night images, uncertainty estimation, real-time deployment.
\end{enumerate}}

% -----------------------------------------------------------------------
% ACKNOWLEDGEMENTS (optional)
% -----------------------------------------------------------------------
\section*{Acknowledgements}
The author thanks Richard Green for project supervision and feedback.
\todo{Add any other acknowledgements (dataset providers, compute resources, etc.)}

% -----------------------------------------------------------------------
% REFERENCES
% -----------------------------------------------------------------------
% Use BibTeX. Add your .bib entries below or in a separate .bib file.
% Recommended citations to track down and add:
%   - ImageNet / ILSVRC paper (Russakovsky et al., 2015)
%   - Your CNN backbone paper (He et al. 2016 for ResNet, etc.)
%   - IM2GPS (Hays & Efros, 2008)
%   - PlaNet (Weyand et al., 2016)
%   - NREL Solar Position Algorithm (Reda & Andreas, 2004) or pysolar
%   - A sky-appearance / sun-direction paper (e.g. Lalonde et al. 2010)
%   - Transfer learning survey (Tan & Yu, 2018 or similar)
\bibliographystyle{IEEEtran}
% \bibliography{references}   % <- uncomment and point to your .bib file

% Inline bibliography for template purposes (replace with BibTeX):
\begin{thebibliography}{00}

\bibitem{placeholder_imagenet}
O.~Russakovsky \textit{et al.}, ``ImageNet Large Scale Visual Recognition
Challenge,'' \textit{Int.\ J.\ Comput.\ Vis.}, vol.~115, no.~3,
pp.~211--252, 2015.

\bibitem{placeholder_resnet}
K.~He, X.~Zhang, S.~Ren, and J.~Sun, ``Deep residual learning for image
recognition,'' in \textit{Proc.\ IEEE CVPR}, pp.~770--778, 2016.

\bibitem{placeholder_im2gps}
J.~Hays and A.~A.~Efros, ``IM2GPS: estimating geographic information from a
single image,'' in \textit{Proc.\ IEEE CVPR}, 2008.

\bibitem{placeholder_planet}
T.~Weyand, I.~Kostrikov, and J.~Philbin, ``PlaNet -- photo geolocation with
convolutional neural networks,'' in \textit{Proc.\ ECCV}, 2016.

\bibitem{placeholder_spa}
I.~Reda and A.~Andreas, ``Solar position algorithm for solar radiation
applications,'' \textit{Solar Energy}, vol.~76, no.~5, pp.~577--589, 2004.

\bibitem{placeholder_lalonde}
J.-F.~Lalonde, A.~A.~Efros, and S.~G.~Narasimhan, ``Estimating the natural
illumination conditions from a single outdoor image,'' \textit{Int.\ J.\
Comput.\ Vis.}, vol.~98, no.~2, pp.~123--145, 2012.

\bibitem{placeholder_forensics}
\todo{Add a relevant image forensics / provenance citation.}

\end{thebibliography}

\end{document}